\documentclass[10pt,a4paper]{amsart}
\usepackage{amsmath,amssymb,amsthm,graphicx}
% Retained minimal preamble for the Stacks Project chapter more-algebra.tex (More on Algebra).
% The chapter file pinned-original-source.tex carries only \input{preamble} on its first line, so no
% preamble text lives inside the chapter itself.  The Stacks Project uses one single shared preamble for
% every chapter; the snapshot stages it identically as sources/stacks-chow/preamble.tex and
% sources/stacks-etale-cohomology/preamble.tex.  The block below is copied mechanically, line for line,
% from sources/stacks-chow/preamble.tex lines 14-18 (the Xy-pic setup) and 154-217 (the native theorem
% declarations together with the mathematical macros).  Nothing was added, removed or redefined; the
% wrapper class and the amsmath/amssymb/amsthm/graphicx base packages are supplied by the assembler.
% Chapter: more-algebra.tex -- More on Algebra.
% For commutative diagrams we use Xy-pic
\usepackage[all]{xy}

% We use 2cell for 2-commutative diagrams.
\xyoption{2cell}

% Theorem environments.
%
\theoremstyle{plain}
\newtheorem{theorem}[subsection]{Theorem}
\newtheorem{proposition}[subsection]{Proposition}
\newtheorem{lemma}[subsection]{Lemma}

\theoremstyle{definition}
\newtheorem{definition}[subsection]{Definition}
\newtheorem{example}[subsection]{Example}
\newtheorem{exercise}[subsection]{Exercise}
\newtheorem{situation}[subsection]{Situation}

\theoremstyle{remark}
\newtheorem{remark}[subsection]{Remark}
\newtheorem{remarks}[subsection]{Remarks}

\numberwithin{equation}{subsection}

% Macros
%
\def\lim{\mathop{\mathrm{lim}}\nolimits}
\def\colim{\mathop{\mathrm{colim}}\nolimits}
\def\Spec{\mathop{\mathrm{Spec}}}
\def\Hom{\mathop{\mathrm{Hom}}\nolimits}
\def\Ext{\mathop{\mathrm{Ext}}\nolimits}
\def\SheafHom{\mathop{\mathcal{H}\!\mathit{om}}\nolimits}
\def\SheafExt{\mathop{\mathcal{E}\!\mathit{xt}}\nolimits}
\def\Sch{\mathit{Sch}}
\def\Mor{\mathop{\mathrm{Mor}}\nolimits}
\def\Ob{\mathop{\mathrm{Ob}}\nolimits}
\def\Sh{\mathop{\mathit{Sh}}\nolimits}
\def\NL{\mathop{N\!L}\nolimits}
\def\CH{\mathop{\mathrm{CH}}\nolimits}
\def\proetale{{pro\text{-}\acute{e}tale}}
\def\etale{{\acute{e}tale}}
\def\QCoh{\mathit{QCoh}}
\def\Ker{\mathop{\mathrm{Ker}}}
\def\Im{\mathop{\mathrm{Im}}}
\def\Coker{\mathop{\mathrm{Coker}}}
\def\Coim{\mathop{\mathrm{Coim}}}

% Boxtimes
%
\DeclareMathSymbol{\boxtimes}{\mathbin}{AMSa}{"02}

%
% Macros for moduli stacks/spaces
%
\def\QCohstack{\mathcal{QC}\!\mathit{oh}}
\def\Cohstack{\mathcal{C}\!\mathit{oh}}
\def\Spacesstack{\mathcal{S}\!\mathit{paces}}
\def\Quotfunctor{\mathrm{Quot}}
\def\Hilbfunctor{\mathrm{Hilb}}
\def\Curvesstack{\mathcal{C}\!\mathit{urves}}
\def\Polarizedstack{\mathcal{P}\!\mathit{olarized}}
\def\Complexesstack{\mathcal{C}\!\mathit{omplexes}}
% \Pic is the operator that assigns to X its picard group, usage \Pic(X)
% \Picardstack_{X/B} denotes the Picard stack of X over B
% \Picardfunctor_{X/B} denotes the Picard functor of X over B
\def\Pic{\mathop{\mathrm{Pic}}\nolimits}
\def\Picardstack{\mathcal{P}\!\mathit{ic}}
\def\Picardfunctor{\mathrm{Pic}}
\def\Deformationcategory{\mathcal{D}\!\mathit{ef}}
\usepackage{cleveref}
\crefname{originallemma}{original lemma}{original lemmas}
\Crefname{originallemma}{Original lemma}{Original lemmas}
\makeatletter
\newlabel{more-algebra-lemma-exact-length-1}{{00MX}{source}{}{sourcepointer-1}{}}
\newlabel{more-algebra-lemma-exact-length-1@cref}{{[originallemma][0][]00MX}{[1][0][]source}{}{}{}}
\makeatother
\crefname{originallemma}{original lemma}{original lemmas}
\Crefname{originallemma}{Original lemma}{Original lemmas}
\makeatletter
\newlabel{algebra-definition-rank}{{00MV}{source}{}{sourcepointer-2}{}}
\newlabel{algebra-definition-rank@cref}{{[originallemma][0][]00MV}{[1][0][]source}{}{}{}}
\makeatother
\crefname{originallemma}{original lemma}{original lemmas}
\Crefname{originallemma}{Original lemma}{Original lemmas}
\makeatletter
\newlabel{algebra-lemma-add-trivial-complex}{{00MT}{source}{}{sourcepointer-3}{}}
\newlabel{algebra-lemma-add-trivial-complex@cref}{{[originallemma][0][]00MT}{[1][0][]source}{}{}{}}
\makeatother
\crefname{originallemma}{original lemma}{original lemmas}
\Crefname{originallemma}{Original lemma}{Original lemmas}
\makeatletter
\newlabel{lemma-auto-ass-implies-P}{{05GN}{source}{}{sourcepointer-4}{}}
\newlabel{lemma-auto-ass-implies-P@cref}{{[originallemma][0][]05GN}{[1][0][]source}{}{}{}}
\makeatother
\crefname{originallemma}{original lemma}{original lemmas}
\Crefname{originallemma}{Original lemma}{Original lemmas}
\makeatletter
\newlabel{lemma-P-universally-injective}{{05GP}{source}{}{sourcepointer-5}{}}
\newlabel{lemma-P-universally-injective@cref}{{[originallemma][0][]05GP}{[1][0][]source}{}{}{}}
\makeatother
\crefname{originallemma}{original lemma}{original lemmas}
\Crefname{originallemma}{Original lemma}{Original lemmas}
\makeatletter
\newlabel{algebra-lemma-weakly-ass}{{0548}{source}{}{sourcepointer-6}{}}
\newlabel{algebra-lemma-weakly-ass@cref}{{[originallemma][0][]0548}{[1][0][]source}{}{}{}}
\makeatother
\crefname{originallemma}{original lemma}{original lemmas}
\Crefname{originallemma}{Original lemma}{Original lemmas}
\makeatletter
\newlabel{algebra-lemma-weakly-ass-zero}{{0588}{source}{}{sourcepointer-7}{}}
\newlabel{algebra-lemma-weakly-ass-zero@cref}{{[originallemma][0][]0588}{[1][0][]source}{}{}{}}
\makeatother
\crefname{originallemma}{original lemma}{original lemmas}
\Crefname{originallemma}{Original lemma}{Original lemmas}
\makeatletter
\newlabel{algebra-lemma-matrix-left-inverse}{{07DQ}{source}{}{sourcepointer-8}{}}
\newlabel{algebra-lemma-matrix-left-inverse@cref}{{[originallemma][0][]07DQ}{[1][0][]source}{}{}{}}
\makeatother
\crefname{originallemma}{original lemma}{original lemmas}
\Crefname{originallemma}{Original lemma}{Original lemmas}
\makeatletter
\newlabel{algebra-proposition-what-exact}{{00N1}{source}{}{sourcepointer-9}{}}
\newlabel{algebra-proposition-what-exact@cref}{{[originallemma][0][]00N1}{[1][0][]source}{}{}{}}
\makeatother
\title{Exactness in length one for a matrix presentation}
\author{The Stacks Project Authors}
\date{stacks-project:more-algebra.tex:local-snapshot-20261012; GNU FDL 1.2 or later}
\begin{document}
\maketitle
\setcounter{subsection}{1}
\noindent\textit{Editorial note.} All delivered mathematical text is the author's own, copied from the pinned Stacks Project chapter file more-algebra.tex; no statement or proof step was rewritten, reordered or removed.\par\medskip
\noindent\textit{Editorial note.} The chapter file pinned-original-source.tex carries only \textbackslash{}input\{preamble\} on its first line, so no preamble text lives inside the chapter. The retained file sources/stacks-more-algebra/retained-preamble.tex is a mechanical, line-for-line copy of the project's single shared preamble as staged in the snapshot (Xy-pic setup plus theorem declarations and mathematical macros). Nothing was added, removed or redefined, and the excerpt serves only as preamble, never as a span coordinate.\par\medskip
\noindent\textit{Editorial note.} Both counted spans start at the \textbackslash{}begin line of their environment, so the statement and the author's proof are delivered whole; surrounding prose of the section is deliberately outside the counted spans and is carried as standard prerequisites instead.\par\medskip
\noindent\textit{Editorial note.} The citation fields (source\_id, source\_version, authors, source\_url, license.name, source\_title) are identical to those of every other card this lane has produced from this chapter, so the batch selector treats them as one source.\par\medskip
\noindent\textit{Allowed standard prerequisites (editorial).}\begin{itemize}
\item Complex attached to a matrix and its exactness in length one. defined in the same section, source lines 3780-3823, before the counted statement. Supplies the object of the statement.
\item A matrix with an invertible maximal minor can be diagonalised by invertible row and column operations. used at proof lines 3843-3875. Is the reduction that makes the verification straightforward.
\item Localisation preserves exactness of a complex of modules. used at proof lines 3875-3896 for the general case. Reduces the statement to the invertible-minor case.
\end{itemize}\medskip
\ifdefined\hypertarget\hypertarget{sourcepointer-1}{}\fi
\noindent\textit{Source reference.} Original lemma 00MX: focused-seeds/AGP/sources/stacks-snapshot-20261011/tags.txt, source lines 839--839. The lemma's own canonical tag, quoted so that readers can resolve the label inside the delivered text; it points at the counted statement itself and carries no proof content.\par\medskip
\ifdefined\hypertarget\hypertarget{sourcepointer-2}{}\fi
\noindent\textit{Source reference.} Original lemma 00MV: focused-seeds/AGP/sources/stacks-snapshot-20261011/tags.txt, source lines 837--837. Navigation pointer to a statement the counted proof refers to but does not contain; nothing load-bearing is delegated to it.\par\medskip
\ifdefined\hypertarget\hypertarget{sourcepointer-3}{}\fi
\noindent\textit{Source reference.} Original lemma 00MT: focused-seeds/AGP/sources/stacks-snapshot-20261011/tags.txt, source lines 835--835. Navigation pointer to a statement the counted proof refers to but does not contain; nothing load-bearing is delegated to it.\par\medskip
\ifdefined\hypertarget\hypertarget{sourcepointer-4}{}\fi
\noindent\textit{Source reference.} Original lemma 05GN: focused-seeds/AGP/sources/stacks-snapshot-20261011/tags.txt, source lines 7067--7067. Navigation pointer to a statement the counted proof refers to but does not contain; nothing load-bearing is delegated to it.\par\medskip
\ifdefined\hypertarget\hypertarget{sourcepointer-5}{}\fi
\noindent\textit{Source reference.} Original lemma 05GP: focused-seeds/AGP/sources/stacks-snapshot-20261011/tags.txt, source lines 7068--7068. Navigation pointer to a statement the counted proof refers to but does not contain; nothing load-bearing is delegated to it.\par\medskip
\ifdefined\hypertarget\hypertarget{sourcepointer-6}{}\fi
\noindent\textit{Source reference.} Original lemma 0548: focused-seeds/AGP/sources/stacks-snapshot-20261011/tags.txt, source lines 6631--6631. Navigation pointer to a statement the counted proof refers to but does not contain; nothing load-bearing is delegated to it.\par\medskip
\ifdefined\hypertarget\hypertarget{sourcepointer-7}{}\fi
\noindent\textit{Source reference.} Original lemma 0588: focused-seeds/AGP/sources/stacks-snapshot-20261011/tags.txt, source lines 6771--6771. Navigation pointer to a statement the counted proof refers to but does not contain; nothing load-bearing is delegated to it.\par\medskip
\ifdefined\hypertarget\hypertarget{sourcepointer-8}{}\fi
\noindent\textit{Source reference.} Original lemma 07DQ: focused-seeds/AGP/sources/stacks-snapshot-20261011/tags.txt, source lines 9420--9420. Navigation pointer to a statement the counted proof refers to but does not contain; nothing load-bearing is delegated to it.\par\medskip
\ifdefined\hypertarget\hypertarget{sourcepointer-9}{}\fi
\noindent\textit{Source reference.} Original lemma 00N1: focused-seeds/AGP/sources/stacks-snapshot-20261011/tags.txt, source lines 843--843. Navigation pointer to a statement the counted proof refers to but does not contain; nothing load-bearing is delegated to it.\par\medskip
\begin{lemma}
\label{lemma-exact-length-1}
Let $(R, \mathfrak m)$ be a local ring.
Suppose that $\varphi : R^m \to R^n$ is a map
of finite free modules. The following are equivalent
\begin{enumerate}
\item $\varphi$ is injective,
\item the rank of $\varphi$ is $m$ and the annihilator of
$I(\varphi)$ in $R$ is zero.
\end{enumerate}
If $R$ is Noetherian these are also equivalent to
\begin{enumerate}
\item[(3)] the rank of $\varphi$ is $m$ and
either $I(\varphi) = R$ or it contains a nonzerodivisor.
\end{enumerate}
Here the rank of $\varphi$ and $I(\varphi)$ are defined
as in Algebra, Definition \ref{algebra-definition-rank}.
\end{lemma}

\begin{proof}
If any matrix coefficient of $\varphi$ is not in $\mathfrak m$, then we apply
Algebra, Lemma \ref{algebra-lemma-add-trivial-complex}
to write $\varphi$ as the sum of $1 : R \to R$ and a map
$\varphi' : R^{m-1} \to R^{n-1}$. It is easy to see that
the lemma for $\varphi'$ implies the lemma for $\varphi$.
Thus we may assume from the outset that all the matrix
coefficients of $\varphi$ are in $\mathfrak m$.

\medskip\noindent
Suppose $\varphi$ is injective. We may assume $m > 0$.
Let $\mathfrak q \in \text{WeakAss}(R)$ so that $R_\mathfrak q$
is an auto-associated ring. Then $\varphi$
induces an injective map $R_\mathfrak q^m \to R_\mathfrak q^n$
which is universally injective by
Lemmas \ref{lemma-auto-ass-implies-P} and \ref{lemma-P-universally-injective}.
Thus $\varphi : \kappa(\mathfrak q)^m \to \kappa(\mathfrak q)^n$
is injective. Hence the rank of $\varphi \bmod \mathfrak q$ is $m$ and
$I(\varphi \otimes \kappa(\mathfrak q))$ is not the zero ideal.
Since $m$ is the maximum rank $\varphi$ can have, we conclude
that $\varphi$ has rank $m$ as well (ranks of matrices
can only drop after base change). Hence
$I(\varphi) \cdot \kappa(\mathfrak q) =
I(\varphi \otimes \kappa(\mathfrak q))$ is not zero.
Thus $I(\varphi)$ is not contained in $\mathfrak q$.
Thus none of the weakly associated primes of $R$
are weakly associated primes of the $R$-module
$\text{Ann}_R I(\varphi)$. Thus $\text{Ann}_R I(\varphi)$
has no weakly associated primes, see
Algebra, Lemma \ref{algebra-lemma-weakly-ass}.
It follows from Algebra, Lemma
\ref{algebra-lemma-weakly-ass-zero}
that $\text{Ann}_R I(\varphi)$ is zero.

\medskip\noindent
Conversely, assume (2). The rank being $m$ implies $n \geq m$.
Write $I(\varphi) = (f_1, \ldots, f_r)$ which is possible as
$I(\varphi)$ is finitely generated. By
Algebra, Lemma \ref{algebra-lemma-matrix-left-inverse}
we can find maps $\psi_i : R^n \to R^m$ such that
$\psi \circ \varphi = f_i \text{id}_{R^m}$. Thus
$\varphi(x) = 0$ implies $f_i x = 0$ for $i = 1, \ldots, r$.
This implies $x = 0$ and hence $\varphi$ is injective.

\medskip\noindent
For the equivalence of (1) and (3) in the Noetherian local case
we refer to Algebra, Proposition \ref{algebra-proposition-what-exact}.
If the ring $R$ is Noetherian but not local, then the reader can
deduce it from the local case; details omitted.
Another option is to redo the argument above using associated
primes, using that there are finitely many of these, using
prime avoidance, and using the characterization of nonzerodivisors
as elements of a Noetherian ring not contained in any associated prime.
\end{proof}

\end{document}
